%! Simplifying Radicals; Adding \& Subtracting — Unit 5, Lesson 5.2 — Guided Notes
\documentclass[10pt]{article}
\usepackage{algebra2-article}
\usepackage{algebra2-boxes}
\newcommand{\TallMath}[1]{$\displaystyle #1\rule[-1.4em]{0pt}{3.2em}$}

\begin{document}

\pageheader{Unit 5, Lesson 5.2}{Guided Notes}
\namedateperiod

\vspace{0.08in}

\begin{objectivebox}
\textbf{By the end of this lesson, I will be able to\dots}
\begin{itemize}\setlength{\itemsep}{2pt}
  \item use the \emph{product} and \emph{quotient} properties of radicals --- and explain why they follow from yesterday's exponent rules rather than being new rules;
  \item write a radical in \emph{simplest form}, with numeric and algebraic radicands and with indices above $2$;
  \item \emph{add and subtract} radicals by collecting \emph{like radicals}, simplifying first so that I can tell which ones are like.
\end{itemize}
\end{objectivebox}

\vspace{0.05in}

\begin{vocabbox}
\small
Fill in each term as we build it together.
\par\vspace{2pt}   % \par is required: \termblanklong's \noindent is a no-op mid-paragraph
\termblanklong{Product Property of Radicals}
\termblanklong{Quotient Property of Radicals}
\termblanklong{Simplest radical form}
\termblanklong{Like radicals}
\end{vocabbox}

\begin{hookbox}
Two of the three statements below are \textbf{true}. One is a lie. Test each one with a decimal
approximation and mark it T or F.
\begin{center}
\renewcommand{\arraystretch}{1.8}
\begin{tabularx}{\linewidth}{@{}X X X@{}}
(a) \; $\sqrt{8}\cdot\sqrt{2} = 4$ & (b) \; $\sqrt{9}+\sqrt{16} = \sqrt{25}$ &
(c) \; $\sqrt{8}+\sqrt{18} = 5\sqrt{2}$ \\
\blank{1.2cm} & \blank{1.2cm} & \blank{1.2cm} \\
\end{tabularx}
\end{center}
The lie is \blank{1.0cm}. Look at what the two true ones have in common and what the false one does
differently: \writeline
\vspace{2pt}
That is today's whole lesson in one line. A radical \textbf{splits over multiplication} --- always.
It \textbf{never} splits over addition. And (c) shows the strange consequence: two radicals that
look completely unrelated can turn out to be the same thing in disguise, but you cannot tell until
you \emph{simplify} them.
\end{hookbox}

\begin{notesbox}{1. Where the Properties Come From (they are yesterday's rules)}
Yesterday: $\sqrt[n]{a}=a^{1/n}$. So watch what happens to a \emph{product} under a radical --- all
we use is the power-of-a-product rule from the Warm-Up:
\[
  \sqrt[n]{ab} \;=\; (ab)^{1/n} \;=\; a^{\,\blank{0.9cm}}\cdot b^{\,\blank{0.9cm}}
  \;=\; \sqrt[n]{a}\cdot\sqrt[n]{b}.
\]
\begin{center}
\begin{tcolorbox}[colback=goldbg, colframe=goldacc, boxrule=0.8pt, arc=1.5mm,
    left=3mm, right=3mm, top=1.5mm, bottom=1.5mm, width=0.90\linewidth]
\centering
\textbf{Product Property} \quad $\displaystyle \sqrt[n]{ab}=\sqrt[n]{a}\cdot\sqrt[n]{b}$
\qquad\qquad
\textbf{Quotient Property} \quad $\displaystyle \sqrt[n]{\frac{a}{b}}=\frac{\sqrt[n]{a}}{\sqrt[n]{b}}$
\;\; $(b\ne 0)$
\end{tcolorbox}
\end{center}
\emph{The fine print (and it matters).} When the index is \textbf{even}, $a$ and $b$ must be
\blank{3.6cm}, or the pieces do not exist as real numbers in the first place. That is Lesson 5.1's
rule doing real work: you may not ``simplify'' $\sqrt{-8}$ into $\sqrt{-4}\cdot\sqrt{2}$, because
\blank{3.0cm} is not a real number.

\vspace{4pt}
\textbf{And there is no third property.} From the Hook:
\begin{center}
\renewcommand{\arraystretch}{1.4}
\begin{tabular}{@{}l l@{}}
$\sqrt{9\cdot 16}=\sqrt{144}=\blank{1.0cm}$ \quad and \quad $\sqrt{9}\cdot\sqrt{16}=\blank{1.0cm}$ &
\quad $\Rightarrow$ \; multiplication: \blank{1.6cm} \\
$\sqrt{9+16}=\sqrt{25}=\blank{1.0cm}$ \quad but \quad $\sqrt{9}+\sqrt{16}=\blank{1.0cm}$ &
\quad $\Rightarrow$ \; addition: \blank{1.6cm} \\
\end{tabular}
\end{center}
\textbf{Write it down and never unlearn it:} \; $\sqrt{a+b}\ne\sqrt{a}+\sqrt{b}$.
\end{notesbox}

\begin{notesbox}{2. Simplest Radical Form, and How to Get There}
A radical is in \textbf{simplest form} when all three of these are true:
\begin{enumerate}[leftmargin=1.7em, itemsep=2pt]
  \item no factor of the radicand is a perfect \blank{1.4cm} power (other than $1$);
  \item there is no \blank{1.6cm} under the radical sign;
  \item there is no radical in a \blank{2.0cm}. \; \emph{(Condition 3 is Lesson 5.3's job.)}
\end{enumerate}

\vspace{3pt}
\textbf{The method: pull out the largest perfect $n$th power.}
\[
  \sqrt{72} \;=\; \sqrt{\blank{1.0cm}\cdot\blank{0.8cm}}
  \;=\; \sqrt{\blank{1.0cm}}\cdot\sqrt{\blank{0.8cm}} \;=\; \blank{1.6cm}
\]
\emph{Why ``largest''?} Suppose you spot only $72=4\cdot 18$ instead. You get
$\sqrt{72}=2\sqrt{18}$ --- which is \textbf{true but not finished}, because $18$ still contains the
perfect square \blank{0.8cm}. Keep going: $2\sqrt{18}=2\cdot\blank{1.2cm}=\blank{1.4cm}$. Same
answer, more steps. \emph{So the last thing you do on every problem is ask: is there still a
perfect square hiding in there?}

\vspace{3pt}
\textbf{Backup method (when you cannot spot it): prime factorization.} Break the radicand all the
way down, then take out one factor for every \textbf{pair} (for a square root):
\[
  \sqrt{72}=\sqrt{2\cdot 2\cdot 2\cdot 3\cdot 3}
  \;\Rightarrow\; \text{a pair of } 2\text{s and a pair of } 3\text{s escape as }
  \blank{0.8cm}\cdot\blank{0.8cm}; \quad \text{one } 2 \text{ is stranded inside.}
\]

\vspace{3pt}
\textbf{Your turn.} Simplify. One of these is already in simplest form --- say which.
\begin{center}
\renewcommand{\arraystretch}{1.8}
\begin{tabularx}{\linewidth}{@{}X X X X@{}}
(a) \; $\sqrt{50}=$ \blank{1.8cm} & (b) \; $\sqrt{48}=$ \blank{1.8cm} &
(c) \; $\sqrt{45}=$ \blank{1.8cm} & (d) \; $\sqrt{30}=$ \blank{1.8cm} \\
(e) \; $\sqrt{200}=$ \blank{1.8cm} & (f) \; $3\sqrt{20}=$ \blank{1.8cm} &
(g) \; $\sqrt{98}=$ \blank{1.8cm} & (h) \; $-5\sqrt{8}=$ \blank{1.8cm} \\
\end{tabularx}
\end{center}
In (f) and (h), the number already outside just \blank{2.6cm} whatever comes out.
\end{notesbox}

\begin{notesbox}{3. Higher Indices: Group in $n$'s, Not in Pairs}
Under a square root you hunt for perfect \emph{squares} and factors escape in \emph{pairs}. Under a
cube root you hunt for perfect \emph{cubes} and factors escape in \blank{1.6cm}. \textbf{The index
tells you the group size.}
\begin{center}
\renewcommand{\arraystretch}{1.3}
\begin{tabularx}{\linewidth}{@{}l X@{}}
Perfect squares & $4,\ 9,\ 16,\ 25,\ 36,\ 49,\ 64,\ 81,\ 100,\ 121,\ 144$ \\
Perfect cubes   & $8,\ 27,\ 64,\ 125,\ 216$ \quad (and $-8,\ -27,\ -64,\dots$) \\
Perfect 4th powers & $16,\ 81,\ 256$ \qquad Perfect 5th powers: $32,\ 243$ \\
\end{tabularx}
\end{center}
\[
  \sqrt[3]{54}=\sqrt[3]{\blank{1.0cm}\cdot\blank{0.8cm}}=\blank{1.6cm}
  \qquad
  \sqrt[3]{16}=\sqrt[3]{\blank{1.0cm}\cdot\blank{0.8cm}}=\blank{1.6cm}
  \qquad
  \sqrt[4]{48}=\blank{1.8cm}
\]
\begin{itemize}[leftmargin=1.5em, itemsep=3pt]
  \item $\sqrt[3]{-24}=\blank{2.0cm}$ \quad --- a negative radicand is fine here because the index
        is \blank{1.2cm} (Lesson 5.1).
  \item \textbf{Trap:} under a cube root, a \emph{pair} of equal factors is worth nothing.
        $\sqrt[3]{4}=\sqrt[3]{2\cdot 2}$ is already in simplest form, because you need
        \blank{1.2cm} twos to get one out.
\end{itemize}
\end{notesbox}

\begin{notesbox}{4. Algebraic Radicands: Divide the Exponent by the Index}
Start with the easy case. $\sqrt{x^{6}}=\blank{1.2cm}$, because $\left(\blank{1.2cm}\right)^2=x^6$.
When the exponent divides evenly by the index, the whole thing comes out.

\vspace{3pt}
When it does not, split off as much as you can:
\[
  \sqrt{x^{7}}=\sqrt{x^{\blank{0.7cm}}\cdot x}=\blank{1.8cm}
\]
\begin{center}
\begin{tcolorbox}[colback=goldbg, colframe=goldacc, boxrule=0.8pt, arc=1.5mm,
    left=3mm, right=3mm, top=1.5mm, bottom=1.5mm, width=0.90\linewidth]
\centering \textbf{Divide the exponent by the index.} The \textbf{quotient} is the exponent that
goes \emph{outside}; the \textbf{remainder} is the exponent that stays \emph{inside}.
\end{tcolorbox}
\end{center}

\vspace{2pt}
\textbf{Why that works --- and why it should feel familiar.} Write the radical as a rational
exponent (Lesson 5.1) and turn the improper fraction into a mixed number, exactly as in Warm-Up
item~3:
\[
  \sqrt{x^{7}} \;=\; x^{7/2} \;=\; x^{\,\blank{0.9cm}} \;=\;
  x^{3}\cdot x^{\,\blank{0.7cm}} \;=\; \blank{1.8cm}
\]
\emph{Simplifying a radical and rewriting an improper fraction as a mixed number are the same
move.} The whole-number part comes out front; the leftover fraction stays under the radical.

\vspace{3pt}
\textbf{Practice the split.}
\begin{center}
\renewcommand{\arraystretch}{1.8}
\begin{tabularx}{\linewidth}{@{}X X X@{}}
(a) \; $\sqrt{y^{9}}=$ \blank{2.0cm} & (b) \; $\sqrt[3]{x^{11}}=$ \blank{2.0cm} &
(c) \; $\sqrt[4]{m^{7}}=$ \blank{2.0cm} \\
\end{tabularx}
\end{center}

\vspace{2pt}
\textbf{Now all of it at once.} Handle the number and each variable separately, then collect.
\[
  \sqrt{50x^{5}y^{8}} \;=\; \sqrt{\blank{1.0cm}\cdot 2}\;\cdot\;\sqrt{x^{4}\cdot x}\;\cdot\;
  \sqrt{y^{8}} \;=\; \blank{3.0cm}
\]
\[
  \sqrt[3]{40a^{7}b^{3}} \;=\; \blank{3.4cm}
\]
\emph{A note we are setting aside:} everything here assumes the variables are \textbf{positive}. If
we did not know that, $\sqrt{x^{2}}$ would have to be written \blank{1.2cm} --- Lesson 5.1's
absolute-value rule.
\end{notesbox}

\begin{notesbox}{5. The Quotient Property}
Same idea, one floor down. Take the root of the top \emph{and} the root of the bottom.
\begin{center}
\renewcommand{\arraystretch}{2.0}
\begin{tabularx}{\linewidth}{@{}X X X@{}}
(a) \; $\sqrt{\dfrac{25}{36}}=$ \blank{1.6cm} &
(b) \; $\sqrt{\dfrac{x^{4}}{9}}=$ \blank{1.6cm} &
(c) \; $\sqrt[3]{\dfrac{8}{27}}=$ \blank{1.6cm} \\
(d) \; $\sqrt{\dfrac{7}{16}}=$ \blank{1.6cm} &
(e) \; $\dfrac{\sqrt{50}}{\sqrt{2}}=$ \blank{1.6cm} &
(f) \; $\dfrac{\sqrt{27}}{\sqrt{3}}=$ \blank{1.6cm} \\
\end{tabularx}
\end{center}
Items (e) and (f) run the property \textbf{backwards} --- two radicals with the same index can be
pulled together under one radical, and often the quotient is much friendlier than either piece.

\vspace{2pt}
\emph{What we cannot finish yet:} $\sqrt{\dfrac{3}{5}}=\dfrac{\sqrt3}{\sqrt5}$ violates simplest-form
condition~3 --- there is a radical in the \blank{2.2cm}. Getting rid of it is called
\textbf{rationalizing}, and it is the first thing we do in Lesson 5.3.
\end{notesbox}

\begin{notesbox}{6. Adding \& Subtracting: Like Radicals}
Look back at Warm-Up item~2. You combined $3x+7x$ into $10x$, but $3x+7y$ went nowhere. Radicals
behave \emph{identically} --- the radical part plays the role of the variable.
\begin{center}
\renewcommand{\arraystretch}{1.5}
\begin{tabular}{@{}l l@{}}
$3x+7x=10x$ & $\qquad 3\sqrt5+7\sqrt5=\blank{1.6cm}$ \\
$3x+7y$ \; cannot combine & $\qquad 3\sqrt5+7\sqrt2$ \; \blank{2.6cm} \\
\end{tabular}
\end{center}
\textbf{Like radicals} have the same \blank{1.6cm} \emph{and} the same \blank{2.0cm}. Add or
subtract the numbers \emph{in front}; the radical itself never changes:
\; $3\sqrt5+7\sqrt5$ is $10\sqrt5$, \textbf{not} $10\sqrt{10}$.

\vspace{4pt}
\textbf{The catch --- and it is the point of putting these two skills in one lesson.} Two radicals
can be like radicals \emph{in disguise}. From the Hook:
\[
  \sqrt{8}+\sqrt{18} \;=\; \blank{1.4cm} + \blank{1.4cm} \;=\; \blank{1.6cm}
\]
Neither term \emph{looked} like the other. \textbf{Simplify first; decide second.} Never answer
``cannot combine'' until both radicals are in simplest form.

\vspace{3pt}
\textbf{Your turn.}
\begin{center}
\renewcommand{\arraystretch}{2.0}
\begin{tabularx}{\linewidth}{@{}X X@{}}
(a) \; $\sqrt{45}-\sqrt{20}=$ \blank{2.4cm} &
(b) \; $5\sqrt{12}+2\sqrt{75}-\sqrt{3}=$ \blank{2.4cm} \\
(c) \; $\sqrt{8x}+\sqrt{50x}=$ \blank{2.4cm} &
(d) \; $\sqrt[3]{16}+\sqrt[3]{54}=$ \blank{2.4cm} \\
(e) \; $\sqrt{7}+\sqrt{28}-\sqrt{2}=$ \blank{2.4cm} &
(f) \; $\sqrt[3]{2}+\sqrt{2}=$ \blank{2.4cm} \\
\end{tabularx}
\end{center}
Item (f) is a warning: even with the same radicand, the \blank{1.4cm} must match too.
\end{notesbox}

\begin{practicebox}
Work these with your class. Assume all variables represent positive real numbers.
\begin{enumerate}[leftmargin=1.7em, itemsep=5pt]
  \item Simplify: \; $\sqrt{108}=\blank{1.8cm}$ \quad $\sqrt[3]{250}=\blank{1.8cm}$ \quad
        $\sqrt[4]{162}=\blank{1.8cm}$
  \item Simplify: \; $\sqrt{75x^{3}}=\blank{2.2cm}$ \quad $\sqrt{12a^{4}b^{7}}=\blank{2.4cm}$ \quad
        $\sqrt[3]{54n^{10}}=\blank{2.4cm}$
  \item Combine: \; $4\sqrt6+9\sqrt6=\blank{1.8cm}$ \quad $\sqrt{27}+\sqrt{48}=\blank{1.8cm}$ \quad
        $3\sqrt{32}-2\sqrt{18}=\blank{1.8cm}$
  \item Use the quotient property: \; $\sqrt{\dfrac{49}{100}}=\blank{1.6cm}$ \quad
        $\dfrac{\sqrt{72}}{\sqrt{2}}=\blank{1.6cm}$
  \item \textbf{Explain.} A classmate simplifies $\sqrt{80}$ as $2\sqrt{20}$ and stops. Are they
        wrong? What would you tell them to check, and what is the finished answer? \writelines{2}
\end{enumerate}
\end{practicebox}

\end{document}
